%% =====================================================================
%%  02805 Social Graphs and Interactions -- Final project paper
%%  Built on the official PNAS Overleaf template (pnas-new.cls).
%%
%%  HOW TO USE
%%  1. In Overleaf open the course template and click "Open as Template":
%%     https://www.overleaf.com/latex/templates/template-for-preparing-your-research-report-submission-to-pnas-using-overleaf/fzcbzjvpvnxn
%%  2. Upload THIS main.tex and references.bib (overwrite main.tex).
%%  3. Make a folder "figures/" and put your figure PDFs/PNGs in it.
%%
%%  HARD RULES (week 8 page)
%%   - Max 6 pages INCLUDING figures and references.
%%   - Sections: Abstract (~150 words), Significance Statement,
%%     Introduction, Results, Discussion, Methods, References,
%%     specific Author Contributions.
%%   - Dataset must contain both a network and text.
%%   - Hand in PDF on DTU Learn by Fri 27 Nov 2026, 23:59.
%%
%%  Placeholders are marked with \todo{...}, which prints in red.
%%  Before hand-in, search for "\todo" and make sure none remain.
%% =====================================================================

% "lineno" prints line numbers: handy for group feedback on drafts.
% Remove it for the version you hand in.
\documentclass[9pt,twocolumn,twoside,lineno]{pnas-new}

\templatetype{pnasresearcharticle}

% ---- Extra packages and helpers --------------------------------------
\usepackage{xcolor}
\usepackage{booktabs}   % nice tables (\toprule, \midrule, \bottomrule)
\newcommand{\todo}[1]{\textcolor{red}{[TODO: #1]}}
\graphicspath{{figures/}}

% ---- Title and authors -----------------------------------------------
% The title states the finding or the question, not the method.
\title{\todo{Title that states the question or the main finding}}

\author[a]{Carlota \todo{Surname}}
\author[a]{\todo{Group member 2}}
\author[a]{\todo{Group member 3}}

\affil[a]{Department of Applied Mathematics and Computer Science (DTU Compute), Technical University of Denmark, 2800 Kgs.\ Lyngby, Denmark}

\leadauthor{\todo{Lead author surname}}

% ---- Significance statement ------------------------------------------
% 2-3 sentences for a reader OUTSIDE the field (max ~120 words).
% Why should anyone care? No jargon, no method names.
\significancestatement{\todo{Two or three plain-language sentences: what question we asked, what we found, and why it matters to someone outside network science.}}

% ---- Author contributions --------------------------------------------
% Must be SPECIFIC. "All members contributed equally" is NOT accepted.
% Name who collected the data, built each analysis, did hand labelling,
% made each figure and wrote which sections.
\authorcontributions{\todo{e.g.\ C.X.\ designed the data collection and collected the network data; A.B.\ built the null model and the community analysis; D.E.\ built the text pipeline and hand-labelled the validation sample; C.X.\ and A.B.\ made Figs.~1--2; D.E.\ made Fig.~3; C.X.\ wrote the Introduction and Methods, A.B.\ the Results, D.E.\ the Discussion; all authors edited the paper.}}
\authordeclaration{The authors declare no competing interests.}
\correspondingauthor{\textsuperscript{1}To whom correspondence should be addressed. E-mail: \todo{sXXXXXX@dtu.dk}}

% 3-5 keywords separated by |
\keywords{\todo{Keyword 1} $|$ \todo{Keyword 2} $|$ network science $|$ text analysis}

% ---- Abstract (WRITE THIS LAST) --------------------------------------
% ~150 words, one paragraph. Must contain:
%   (1) the question, (2) the data (source + size),
%   (3) the main result WITH A NUMBER, (4) the check behind it,
%   (5) what it means.
% The one-sentence finding here must match the Results and Discussion.
\begin{abstract}
\todo{Question in one sentence.} \todo{Data: what, from where, how many nodes, edges and documents.} \todo{Main result with a number, e.g.\ ``X is 2.3 times higher than in a degree-preserving null (95\% interval 1.8--2.9)''.} \todo{The check that makes us believe it.} \todo{What it means, and its main limitation.}
\end{abstract}

\dates{This manuscript was compiled on \today}
\doi{02805 Social Graphs and Interactions, DTU, Fall 2026}

\begin{document}

\maketitle
\thispagestyle{firststyle}
\ifthenelse{\boolean{shortarticle}}{\ifthenelse{\boolean{singlecolumn}}{\abscontentformatted}{\abscontent}}{}

%% =====================================================================
\section*{Introduction}
%% =====================================================================
% Rough budget: ~500 words. Four moves, in this order:
%  1. The question, and why it matters.
%  2. What is already known: CITE the work you build on
%     (graded under "References and literature").
%  3. The gap: what is still unknown, to you AND to the domain's experts.
%  4. One-sentence preview of the answer.

\dropcap{\todo{F}}irst paragraph: \todo{the question and why it matters.}

\todo{What is already known, with citations, e.g.\ prior work on the topic \cite{wagner2015man,ribeiro2020auditing}. Cite the methods you rely on as you introduce them \cite{coscia2025atlas}.}

\todo{The gap: what this prior work does not answer.}

\todo{Here we \ldots\ We find that \ldots\ (one-sentence preview, consistent with the abstract).}

%% =====================================================================
\section*{Results}
%% =====================================================================
% Rough budget: ~1200 words + figures. The FIGURES carry the Results.
%  - "Three good figures beat eight."
%  - Each figure answers part of the question and shows the CHECK beside
%    the claim (effect vs.\ its null, a second method, or hand labels).
%  - A 1000-node network drawing is decoration; effect-vs-null is a result.
%  - Use subsections whose titles are FINDINGS, not method names.

\subsection*{\todo{Finding 1 stated as a sentence (network half)}}
\todo{What we measured, what it captures and what it misses, the result with a number, and how it compares with the null (Fig.~\ref{fig:main}).}

\begin{figure*}[t]
\centering
% \includegraphics[width=0.95\linewidth]{fig1_main}
\fbox{\parbox[c][5cm][c]{0.9\linewidth}{\centering\todo{Main figure: the effect next to its null / check}}}
\caption{\textbf{\todo{One-line takeaway of the figure.}} (A) \todo{What is plotted, axes and units.} (B) \todo{The comparison: the null model or second method, e.g.\ shaded band = 95\% range over 1000 degree-preserving rewirings.} \todo{What the reader should see.}}
\label{fig:main}
\end{figure*}

\subsection*{\todo{Finding 2 stated as a sentence (text half)}}
\todo{Text measurement, its validation against hand labels or a second method (Fig.~\ref{fig:text}), and the result with a number.}

\begin{figure}[t]
\centering
% \includegraphics[width=\linewidth]{fig2_text}
\fbox{\parbox[c][4cm][c]{0.9\linewidth}{\centering\todo{Figure 2: text result}}}
\caption{\textbf{\todo{Takeaway.}} \todo{What is plotted, the comparison, and what to see.}}
\label{fig:text}
\end{figure}

\subsection*{\todo{Finding 3: how the network and text halves connect}}
\todo{The result that only exists because the network and the text are combined. This is what makes the project more than two mini-projects.}

\begin{figure}[t]
\centering
% \includegraphics[width=\linewidth]{fig3_combined}
\fbox{\parbox[c][4cm][c]{0.9\linewidth}{\centering\todo{Figure 3: network $\times$ text}}}
\caption{\textbf{\todo{Takeaway.}} \todo{What is plotted, the comparison, and what to see.}}
\label{fig:combined}
\end{figure}

% Optional: a compact table of the data or of the effect sizes against nulls.
\begin{table}[t]
\centering
\caption{\todo{Dataset summary, or effect sizes against their nulls.}}
\label{tab:data}
\begin{tabular}{lrr}
\toprule
 & \todo{Group A} & \todo{Group B} \\
\midrule
Nodes & -- & -- \\
Edges & -- & -- \\
Documents (texts) & -- & -- \\
Tokens & -- & -- \\
\bottomrule
\end{tabular}
\end{table}

%% =====================================================================
\section*{Discussion}
%% =====================================================================
% Rough budget: ~400 words.
%  - What the result DOES mean, and what it does NOT mean.
%  - At least one limitation you found YOURSELF, not a generic one.
%  - Relate the result back to the prior work cited in the Introduction.
%  - A negative result with a convincing check is a fine result.

\todo{Main finding restated in one sentence, the same as in the abstract.}

\todo{What it does not mean, e.g.\ correlation vs.\ causation, or which population the sample represents.}

\todo{Limitation we found ourselves, and how large it could be.}

\todo{Back to the literature, and what the next study would be.}

%% =====================================================================
\section*{Materials and Methods}
%% =====================================================================
% Rough budget: ~700 words. Enough for a reader to REDO the work.

\subsection*{Data collection}
\todo{Source and API, collection date(s), query and seed list, pauses and rate limits, and what was dropped and why. ``Keep what you collected and when.''}

\subsection*{Network construction}
\todo{What a node is, what an edge is, directed or weighted or projected, filtering thresholds, and final size (Table~\ref{tab:data}).}

\subsection*{Text processing}
\todo{Tokenisation, stop words, lemmatisation, language filter, and the dictionary, model or method used for each text measure.}

\subsection*{Measurements}
\todo{For each measure: what it is, what it captures, and what it misses.}

\subsection*{Validation and null models}
\todo{For each claim, its check and the result of that check: null model (e.g.\ degree-preserving rewiring \cite{maslov2002specificity}, number of realisations), second method, and hand-labelled sample (size, who labelled, agreement).}

\subsection*{Code and data availability}
\todo{Link to the group repository / website with code and the processed data.}

%% ---- References -------------------------------------------------------
% Cite what you used: data sources, tools (NetworkX, etc.), methods and
% prior work. Real, relevant and correctly cited (graded).
\bibliography{references}

\end{document}
